\documentclass{subfiles}
\begin{document}
    \begin{procedure}
        \caption{RSA()}\label{proc:RSA-schema}
        \nl $\B$ chooses two prime numbers $p \text{ and } q$ with similar size,
            and computes $n = p * q$ and $\lambda = \lcm(p - 1, q - 1)$.
            Additionally $\B$ chooses $\epsilon$ such that $\gcd(\lambda, \epsilon) = 1$.
            The pair $(n, \epsilon)$ represents $\B$'s public key.
            \DontPrintSemicolon\;

        \nl $\A$ cyphers its message $m$, which is assumed to be an integer,
            by computing $c = m^{\epsilon} \mod[n]$, and sends it to $\B$.
            \DontPrintSemicolon\;

        \nl To retrieve $m$, $\B$ computes $m = c^{d} \mod[n]$,
            where $d$ is the Bézout coeffient of $\epsilon$ in the identity $1 = d\epsilon + \mu\lambda$.
            \DontPrintSemicolon\;

        \nl If $\E$ gets to $c$, although has access to $\B$'s public key,
            can't possibly retrieve $m$. The only possibility is for $\E$ to factor $n$,
            which is a well known hard problem.
    \end{procedure}
\end{document}
